% Generated draft: requires mathematical and source review.
\documentclass[11pt]{article}
\usepackage[margin=25mm]{geometry}
\usepackage{fontspec}
\setmainfont{Latin Modern Roman}
\usepackage{amsmath,amssymb,mathtools,unicode-math}
\setmathfont{Latin Modern Math}
\usepackage{enumitem,needspace,xurl,hyperref}
\setlength{\parindent}{0pt}
\setlength{\parskip}{5pt}
\setlength{\emergencystretch}{3em}
\newcommand{\N}{\mathbb{N}}
\newcommand{\Z}{\mathbb{Z}}
\newcommand{\Q}{\mathbb{Q}}
\newcommand{\R}{\mathbb{R}}
\newcommand{\C}{\mathbb{C}}
\newcommand{\F}{\mathbb{F}}
\newcommand{\A}{\mathbb{A}}
\newcommand{\PP}{\mathbb{P}}
\newcommand{\E}{\mathbb{E}}
\newcommand{\eps}{\varepsilon}
\newcommand{\dd}{\,\mathrm{d}}
\begin{document}
\section*{UnsolvedMath ID: 800084}
\textbf{Source problem code:} catalogue entry
\subsection{Definitions and mathematical statement}
Let \(p\colon S^7\to S^4\) be a given ``generalized'' Hopf fibration with fibre \(S^3\), as in the supplied source. The source does not define ``generalized'', specify the manifold category, or identify the structure group of this fibration; these qualifications remain unresolved.

Under the usual integral-homology interpretation, a homology \(4\)-sphere is a closed, connected, oriented \(4\)-manifold \(H\) satisfying
\[
H_i(H;\Z)\cong H_i(S^4;\Z)\qquad\text{for every }i\geq 0.
\]
For chosen orientations, a continuous map \(f\colon H\to S^4\) has degree one if
\[
f_*[H]=[S^4]\quad\text{in }H_4(S^4;\Z).
\]
The integral coefficients and manifold conventions in this interpretation are not explicitly specified in the source; any additional regularity requirement on \(H\) or \(f\) must also be clarified.

The pullback of \(p\) along \(f\) is the projection
\[
f^*S^7=\{(h,x)\in H\times S^7\mid f(h)=p(x)\}\longrightarrow H,
\qquad (h,x)\longmapsto h.
\]
In the usual bundle-theoretic sense, this \(S^3\)-bundle is flat if it admits local trivializations whose transition functions are locally constant maps into the intended group of automorphisms of \(S^3\). The source does not specify that group or whether these automorphisms must be homeomorphisms, diffeomorphisms, or satisfy additional conditions.

\textbf{Question (Freedman, as attributed in the supplied record).} For every given generalized Hopf fibration
\[
S^3\longrightarrow S^7\xrightarrow{\,p\,}S^4,
\]
does there exist a homology \(4\)-sphere \(H\) and a degree-one map \(f\colon H\to S^4\) such that \(f^*S^7\to H\) is a flat \(S^3\)-bundle? This is a conditional formulation: a fully specified mathematical statement requires clarification of ``generalized'', the category and structure group, and the homology convention.
\subsection{Short English statement}
Given a generalized Hopf fibration \(S^3\to S^7\to S^4\), can one find a homology \(4\)-sphere mapping to \(S^4\) with degree one so that the pulled-back \(S^3\)-bundle is flat? The source leaves the meaning of ``generalized'' and the relevant category and structure group unspecified. Integral homology is used above as an explicitly qualified conventional interpretation.
\subsection{Sources}
\begin{itemize}
\item Supplied local UnsolvedMath registry record, permanent ID \(800084\), exact source problem code \texttt{catalogue entry}, titled ``Foliations and Diffeomorphism Groups --- Question 25 (Freedman)''. The record lists Nariman and Yazdi as proposers, with proposed year 2025, and gives attribution/context to Freedman.
\item Nariman and Yazdi, \emph{Problem List on Foliations and Diffeomorphism Groups} (2025), Question 25. Supplied source URL: \url{https://arxiv.org/abs/2503.18274}. The record reports access on 2026-10-16 and extraction by \texttt{arxiv\_latex}; these are supplied provenance statements, not claims of independent source access or verification.
\item The supplied registry labels this problem \texttt{open}. This is a supplied status label, not an independently verified assessment of the literature. No separate literature references or dated literature assessment are supplied.
\end{itemize}
\end{document}
